Táto práca sa zaoberala návrhom a vývojom mobilnej aplikácie Joinly, ktorej cieľom bolo uľahčiť objavovanie a správu spoločenských udalostí s využitím moderných technológií pre platformu Android.                                                                                                                                                                                                                                                                   
  V úvodnej kapitole boli vysvetlené základné pojmy potrebné na pochopenie problematiky. Následne bol analyzovaný súčasný stav — existujúce riešenia a ich nedostatky, ktoré motivovali vznik tejto aplikácie. Ďalšia kapitola sa   
  venovala použitým technológiám a nástrojom. Na základe týchto poznatkov boli definované funkčné a nefunkčné požiadavky, ktoré tvorili základ pre návrh riešenia. Návrh bol zdokumentovaný prostredníctvom UML diagramov. Následne 
  prebehla samotná implementácia aplikácie a v záverečnej kapitole bolo vykonané testovanie, ktoré overilo správnosť implementovaných funkcionalít.

  Výsledkom práce je funkčná aplikácia pokrývajúca základné požiadavky — autentifikáciu používateľov, správu udalostí, systém účasti, odporúčací systém, chat aj moderačné nástroje. Napriek tomu aplikácia v súčasnom stave nie je 
  plne pripravená na masové nasadenie. Využíva bezplatné úrovne troch cloudových poskytovateľov — Firebase, Google Maps a Cloudinary — ktoré sú vhodné pre vývoj a testovanie, no pri väčšom počte používateľov by narazili na      
  limity výkonu aj nákladov. Do budúcnosti by bolo vhodné pre jednotlivé funkcie využiť špecializovaných poskytovateľov optimalizovaných na škálovateľnosť, prípadne vyvinúť vlastné riešenia prispôsobené konkrétnym potrebám      
  aplikácie.

  Ďalším obmedzením je výpočet odporúčacieho systému, ktorý prebieha priamo na zariadení používateľa. Pri väčšom množstve dát by bolo vhodnejšie presunúť túto logiku na server formou cloudových funkcií. Aplikácia taktiež prešla 
  len obmedzeným testovaním v kontrolovanom prostredí — dlhodobé testovanie reálnymi používateľmi by odhalilo ďalšie nedostatky, najmä v oblasti odporúčacieho systému, ktorý bez dostatočného množstva reálnych interakcií nie je  
  možné plnohodnotne vyhodnotiť. Do budúcnosti je možné aplikáciu rozšíriť o prémiové služby pre organizátorov, nákup vstupeniek priamo v aplikácii, ďalšie spôsoby overenia identity, ako aj o riešenie legislatívnych požiadaviek 
  v oblasti GDPR a ochrany osobných údajov.

  Napriek týmto obmedzeniam predstavuje aplikácia Joinly funkčný základ, ktorý demonštruje reálnu využiteľnosť moderných mobilných technológií pri riešení každodenných potrieb. Aplikácia prináša hodnotu najmä v spojení komunít  
  okolo spoločných záujmov, zjednodušení organizácie udalostí a personalizovanom objavovaní podujatí prostredníctvom odporúčacieho systému. Týmto spôsobom napĺňa svoj hlavný cieľ — priblížiť ľudí k udalostiam, ktoré sú pre nich 
  relevantné.
